% ===============================================================
% GRADE 10 -- MATATAG Science 10, Three-Term BOW (rev. 10 Apr 2026)
% plus appendix of K-12 MELC-only topics (see sources.tex)
% ===============================================================
\section{Grade 10}
\resetcards
{\small\textit{Scope note:} DepEd's original phase-in schedule puts Grade 10 MATATAG in SY 2027--2028, but a Grade 10 MATATAG Three-Term BOW exists for SY 2026--2027. The core below follows that BOW. Topics found only in the older K-12 MELCs are in the appendix at the end of this section. \verify{Confirm with your school which Grade 10 curriculum applies this year.}\par}
\begin{gradebody}

% ---------------------------------------------------------------
\subsection{Term 1: Chemical Reactions, Homeostasis, Evolution}

\topic{Evidence of Reactions; Acids, Bases, and Indicators}
\kf
\begin{keyfacts}
\item Signs of a chemical reaction: colour change, precipitate, gas (bubbles), odour, temperature change (and sometimes light).
\item Acids release \ce{H+} and taste sour. Bases release \ce{OH-}, taste bitter, and feel slippery. Acid + base $\rightarrow$ salt + water (\textbf{neutralization}).
\item Common acids: \ce{HCl} (stomach acid, muriatic acid), acetic acid (vinegar/\emph{suka}), citric acid (calamansi). Common bases: \ce{NaOH} (lye, caustic soda), \ce{Mg(OH)2} (antacid), ammonia (cleaners). Salt solution: saline (\ce{NaCl}).
\item Indicators: phenolphthalein is colourless in acid and pink in base. Methyl orange is red in acid and yellow in base. Red cabbage is red in acid and green-yellow in base. Turmeric (\emph{luyang dilaw}) is yellow, turning reddish-brown in base.
\end{keyfacts}
\cards
\qa{Five signs of a chemical reaction}{Colour, gas, precipitate, odour, temperature change}{A}
\qa{Insoluble solid formed from two solutions}{Precipitate}{E}
\qa{Phenolphthalein in a base}{Pink}{A}
\qa{Phenolphthalein in an acid}{Colourless}{A}
\qa{Acid in the stomach}{Hydrochloric acid}{E}
\qa{Acid in vinegar}{Acetic acid}{E}
\qa{Common name of \ce{NaOH}}{Lye (caustic soda)}{A}
\qa{Chemical name of baking soda}{Sodium bicarbonate (\ce{NaHCO3})}{A}
\qa{Acid + base gives}{Salt + water}{E}
\qa{Base in milk of magnesia antacid}{Magnesium hydroxide}{A}
\qa{Turmeric indicator in a base turns}{Reddish-brown}{D}

\topic{Types of Chemical Reactions}
\kf
\begin{keyfacts}
\item \textbf{Combination (synthesis)}: $A + B \rightarrow AB$, e.g.\ \ce{2Mg + O2 -> 2MgO}.
\item \textbf{Decomposition}: $AB \rightarrow A + B$, e.g.\ \ce{2H2O2 -> 2H2O + O2}.
\item \textbf{Single replacement}: $A + BC \rightarrow AC + B$, e.g.\ \ce{Zn + 2HCl -> ZnCl2 + H2} (acid + metal gives hydrogen).
\item \textbf{Double replacement}: $AB + CD \rightarrow AD + CB$, e.g.\ \ce{AgNO3 + NaCl -> AgCl v + NaNO3}.
\item \textbf{Combustion}: fuel + \ce{O2} $\rightarrow$ \ce{CO2} + \ce{H2O}, e.g.\ \ce{CH4 + 2O2 -> CO2 + 2H2O}.
\item Acid + carbonate gives salt + water + \ce{CO2} (vinegar + baking soda).
\item Gas tests: \ce{H2} gives a ``squeaky pop'' with a burning splint; \ce{O2} relights a glowing splint.
\item Rusting needs iron, oxygen, and water. Galvanizing (a zinc coating) prevents it.
\end{keyfacts}
\cards
\qa{$A + B \rightarrow AB$}{Combination (synthesis)}{E}
\qa{$AB \rightarrow A + B$}{Decomposition}{E}
\qa{$A + BC \rightarrow AC + B$}{Single replacement}{E}
\qa{$AB + CD \rightarrow AD + CB$}{Double replacement}{E}
\qa{Burning LPG is what type of reaction?}{Combustion}{E}
\qa{Gas from acid + metal}{Hydrogen}{A}
\qa{Gas from vinegar + baking soda}{Carbon dioxide}{E}
\qa{Products of complete hydrocarbon combustion}{\ce{CO2} + water}{A}
\qa{Test for hydrogen}{Squeaky pop with a burning splint}{D}
\qa{Test for oxygen}{Relights a glowing splint}{D}
\qa{Three things needed for rusting}{Iron, oxygen, water}{A}
\qa{Zinc coating that protects iron}{Galvanizing}{D}
\begin{confbox}
\confusion{Single}{double replacement}{single: an element swaps into a compound. Double: two compounds swap partners (often making a precipitate).}
\confusion{Physical}{chemical change}{physical changes form or state only (melting, dissolving); chemical changes make a new substance.}
\end{confbox}

\topic{Chemical Equations and Conservation of Mass}
\kf
\begin{keyfacts}
\item \textbf{Antoine Lavoisier}: mass is neither created nor destroyed in a reaction. Atoms only rearrange.
\item Write a word equation, then a formula equation, then balance it by changing \textbf{coefficients only}, never subscripts.
\item \ce{2H2 + O2 -> 2H2O}; \ce{N2 + 3H2 -> 2NH3}.
\end{keyfacts}
\cards
\qa{Law of conservation of mass}{Antoine Lavoisier}{E}
\qa{What you may change when balancing}{Coefficients only}{A}
\qa{Balance \ce{H2 + O2 -> H2O}}{2, 1, 2}{A}
\qa{Balance \ce{N2 + H2 -> NH3}}{1, 3, 2}{D}
\qa{Sealed flask reaction: mass after vs.\ before}{The same}{E}
\begin{confbox}
\confusion{Coefficient}{subscript}{a coefficient (\ce{2H2O}) counts molecules; a subscript (\ce{H2O}) is part of the identity, so changing it changes the substance.}
\end{confbox}

\topic{Rates of Reaction; Exothermic and Endothermic}
\kf
\begin{keyfacts}
\item \textbf{Collision theory}: particles must collide with enough energy (the activation energy) and the right orientation.
\item Rate factors: temperature, concentration, surface area (particle size), and catalysts.
\item A \textbf{catalyst} speeds a reaction by lowering the activation energy and is not used up. Enzymes are biological catalysts.
\item Applications: refrigeration slows spoilage, the fire triangle (heat, fuel, oxygen), and painting or galvanizing slows corrosion.
\item \textbf{Exothermic} reactions release heat (combustion, respiration, neutralization). \textbf{Endothermic} reactions absorb heat (photosynthesis, cold packs).
\end{keyfacts}
\cards
\qa{Four factors affecting reaction rate}{Temperature, concentration, surface area, catalyst}{A}
\qa{Why we refrigerate food}{Low temperature slows reactions}{E}
\qa{Speeds a reaction without being used up}{Catalyst}{E}
\qa{Biological catalysts}{Enzymes}{E}
\qa{A catalyst works by lowering the}{Activation energy}{D}
\qa{Why powdered reactants react faster than lumps}{Greater surface area}{A}
\qa{Reaction that releases heat}{Exothermic}{E}
\qa{Cold pack or photosynthesis}{Endothermic}{A}
\qa{Three sides of the fire triangle}{Heat, fuel, oxygen}{E}
\qa{Particles must collide with enough energy and correct orientation}{Collision theory}{D}
\begin{confbox}
\confusion{Exothermic}{endothermic}{exo means heat \emph{exits}, so the surroundings get warmer; endo means heat \emph{enters}, so the surroundings get colder.}
\end{confbox}

\topic{Homeostasis and Feedback}
\kf
\begin{keyfacts}
\item \textbf{Homeostasis} (term by Walter Cannon) is a stable internal balance. Indicators include body temperature (about \SI{37}{\celsius}), blood glucose, and blood pressure (about 120/80\,mmHg).
\item A feedback loop runs receptor $\rightarrow$ control centre (brain/hypothalamus) $\rightarrow$ effector.
\item \textbf{Negative feedback} reverses a change and is the most common type: thermoregulation, and blood sugar control by \textbf{insulin} (lowers) and \textbf{glucagon} (raises), both from the pancreas.
\item \textbf{Positive feedback} amplifies a change: childbirth contractions (oxytocin), blood clotting.
\item When hot: sweating and vasodilation. When cold: shivering and vasoconstriction.
\end{keyfacts}
\cards
\qa{Maintaining a stable internal environment}{Homeostasis}{E}
\qa{Coined ``homeostasis''}{Walter Cannon}{D}
\qa{Normal body temperature}{About \SI{37}{\celsius}}{E}
\qa{Brain part that acts as the body's thermostat}{Hypothalamus}{A}
\qa{Hormone that lowers blood sugar}{Insulin}{E}
\qa{Hormone that raises blood sugar}{Glucagon}{A}
\qa{Organ that secretes insulin}{Pancreas}{E}
\qa{Feedback that reverses a change}{Negative feedback}{E}
\qa{Example of positive feedback}{Childbirth contractions}{A}
\qa{Blood clotting uses which feedback?}{Positive}{D}
\qa{Body's responses to overheating}{Sweating, vasodilation}{A}
\qa{Shivering is a response to}{Cold}{E}
\qa{Part of the loop that detects change}{Receptor}{A}
\qa{Part of the loop that carries out the response}{Effector}{D}
\qa{Typical normal adult blood pressure}{120/80\,mmHg}{A}
\qa{Disease of poor blood-sugar control}{Diabetes mellitus}{E}
\begin{confbox}
\confusion{Negative}{positive feedback}{negative brings things back to normal; positive pushes further from normal until an endpoint.}
\confusion{Insulin}{glucagon}{\textbf{ins}ulin puts sugar \textbf{in}to cells (lowers it); glucagon releases stored glucose (raises it).}
\confusion{Vasodilation}{vasoconstriction}{dilation widens vessels to lose heat; constriction narrows them to keep heat.}
\end{confbox}

\topic{Mechanisms and Evidence of Evolution}
\kf
\begin{keyfacts}
\item \textbf{Charles Darwin}: natural selection; \emph{On the Origin of Species} (1859); HMS \emph{Beagle}; Gal\'apagos finches. \textbf{Alfred Russel Wallace} reached the same idea independently.
\item \textbf{Lamarck}: use and disuse, inheritance of acquired traits (now rejected).
\item Natural selection: overproduction, variation, competition, survival and reproduction of the fittest, then inheritance.
\item Key concepts: variation, heredity, isolation (leading to new species), selection, adaptation.
\item Evidence: fossils (transitional \emph{Archaeopteryx}), biogeography, and comparative morphology: homologous structures (same origin), analogous structures (same function), vestigial structures (appendix, tailbone).
\item Classic examples: the peppered moth and antibiotic resistance.
\end{keyfacts}
\cards
\qa{Primary mechanism of evolution}{Natural selection}{E}
\qa{Darwin's 1859 book}{\emph{On the Origin of Species}}{E}
\qa{Ship Darwin sailed on}{HMS \emph{Beagle}}{A}
\qa{Islands of Darwin's finches}{Gal\'apagos}{E}
\qa{Co-discoverer of natural selection}{Alfred Russel Wallace}{A}
\qa{Theory of use and disuse}{Lamarck}{E}
\qa{Coined ``survival of the fittest''}{Herbert Spencer}{D}
\qa{Same structure, different function (human arm, bat wing)}{Homologous}{E}
\qa{Same function, different origin (bird vs.\ insect wing)}{Analogous}{A}
\qa{Reduced, nearly functionless organ (appendix)}{Vestigial}{E}
\qa{Reptile--bird transitional fossil}{\emph{Archaeopteryx}}{A}
\qa{Study of where species are distributed}{Biogeography}{A}
\qa{Bacteria surviving antibiotics is an example of}{Natural selection}{A}
\qa{Moth that darkened in industrial England}{Peppered moth}{D}
\qa{Separation of populations leading to new species}{Isolation $\rightarrow$ speciation}{D}
\qa{Inherited trait that improves survival}{Adaptation}{E}
\qa{``Fittest'' means}{Most surviving offspring}{D}
\begin{confbox}
\confusion{Darwin}{Lamarck}{Darwin: individuals with existing variations survive. Lamarck: traits gained during life are passed on (wrong).}
\confusion{Homologous}{analogous}{homologous means same ancestry; analogous means same job but different ancestry.}
\confusion{Evolution}{natural selection}{evolution is the change over generations; natural selection is the main \emph{mechanism}.}
\end{confbox}

% ---------------------------------------------------------------
\subsection{Term 2: Populations, Biotechnology, Climate}

\topic{Carrying Capacity and Population Growth}
\kf
\begin{keyfacts}
\item \textbf{Carrying capacity} ($K$) is the largest population an environment can sustain.
\item Limiting factors include food, water, shelter, mates, and space.
\item Density-dependent factors include competition, disease, and predation. Density-independent factors include typhoons, fires, and drought.
\item Exponential growth gives a J-curve. Logistic growth gives an S-curve that levels off at $K$, so growth slows as the population nears $K$.
\item Population changes through births, deaths, immigration (in), and emigration (out).
\end{keyfacts}
\cards
\qa{Largest population an environment can support}{Carrying capacity}{E}
\qa{Symbol for carrying capacity}{$K$}{A}
\qa{Growth with unlimited resources}{J-curve (exponential)}{A}
\qa{Growth that levels off at $K$}{S-curve (logistic)}{A}
\qa{A factor such as food that restricts growth}{Limiting factor}{E}
\qa{Near $K$, population growth}{Slows down}{E}
\qa{Moving into a population}{Immigration}{E}
\qa{A typhoon kills animals regardless of crowding}{Density-independent}{D}
\qa{Disease spreads faster in crowded populations}{Density-dependent}{D}
\qa{Purpose of a closed fishing season}{Let fish stocks recover}{A}
\begin{confbox}
\confusion{Immigration}{emigration}{immigration means \textbf{i}n; emigration means \textbf{e}xit.}
\confusion{Density-dependent}{independent}{dependent factors get stronger with crowding; independent factors hit regardless of crowding.}
\end{confbox}

\topic{Biotechnology}
\kf
\begin{keyfacts}
\item \textbf{Biotechnology} uses living organisms or their processes to make products.
\item Traditional biotechnology is fermentation: bread (yeast), vinegar/\emph{suka}, soy sauce/\emph{toyo}, \emph{patis}, \emph{bagoong}, \emph{kesong puti}, \emph{basi}, and nata de coco (bacterial cellulose).
\item Modern biotechnology uses genetic engineering (GMOs), for example Golden Rice (beta-carotene), Bt corn (an insect-resistance gene from \emph{Bacillus thuringiensis}), and human insulin made by bacteria.
\item Reproductive technology and cloning: IVF (Louise Brown, 1978, first ``test-tube baby''); Dolly the sheep (1996, first mammal cloned from an adult cell). CRISPR is a gene-editing tool.
\item PH research: IRRI and PhilRice. The ethical debate covers safety, ecology, access, and religion.
\end{keyfacts}
\cards
\qa{Microbe in bread and alcohol fermentation}{Yeast}{E}
\qa{Fermented coconut-water gel}{Nata de coco}{E}
\qa{Bacterium that makes nata de coco}{\emph{Acetobacter} (\emph{Komagataeibacter}) \emph{xylinus}}{D}
\qa{Fresh white cheese from carabao's milk}{Kesong puti}{A}
\qa{Ilocano sugarcane wine}{Basi}{A}
\qa{Patis is traditional or modern biotechnology?}{Traditional (fermentation)}{E}
\qa{Rice engineered to make beta-carotene}{Golden Rice}{E}
\qa{``Bt'' in Bt corn stands for}{\emph{Bacillus thuringiensis}}{D}
\qa{Bt corn resists}{Insect pests (corn borer)}{A}
\qa{Organism with deliberately altered DNA}{GMO}{E}
\qa{Human insulin is mass-produced using}{Genetically engineered bacteria}{A}
\qa{First mammal cloned from an adult cell}{Dolly the sheep}{A}
\qa{First IVF baby (1978)}{Louise Brown}{D}
\qa{Gene-editing tool (Nobel Prize 2020)}{CRISPR--Cas9}{D}
\qa{International rice institute in Los Ba\~nos}{IRRI}{A}
\begin{confbox}
\confusion{Traditional}{modern biotech}{traditional uses whole organisms naturally (fermentation); modern changes DNA directly.}
\confusion{GMO}{hybrid}{a hybrid comes from crossbreeding (selective breeding); a GMO has genes inserted in a lab.}
\confusion{Cloning}{IVF}{cloning copies one parent's DNA; IVF is fertilization outside the body, with two parents.}
\end{confbox}

\topic{Global Climate, ENSO, and Sustainability}
\kf
\begin{keyfacts}
\item The \textbf{greenhouse effect} is natural. Extra greenhouse gases (\ce{CO2}, \ce{CH4}, \ce{N2O}, water vapour, fluorinated gases) enhance it and cause global warming.
\item Evidence: rising temperatures, sea-level rise (thermal expansion and melting land ice), shrinking glaciers, ocean acidification, and stronger extremes. The Keeling Curve tracks rising \ce{CO2} at Mauna Loa.
\item Global warming so far is \verify{about \SI{1.1}{\celsius} above pre-industrial levels (IPCC AR6, 2011--2020)}.
\item Agreements: Kyoto Protocol (1997) and the Paris Agreement (2015; well below \SI{2}{\celsius}, aiming for \SI{1.5}{\celsius}). The Montreal Protocol bans CFCs (ozone).
\item \textbf{ENSO}: El Ni\~no is a warm central/eastern Pacific, bringing below-normal rain and drought to the PH. La Ni\~na is the cool phase, bringing above-normal rain.
\item PH laws: Climate Change Act (RA 9729) and Renewable Energy Act (RA 9513). Renewables: geothermal, hydro, solar, wind (Bangui, Ilocos Norte), biomass.
\item \textbf{Mitigation} reduces emissions; \textbf{adaptation} adjusts to the impacts.
\end{keyfacts}
\cards
\qa{Main human-made greenhouse gas}{Carbon dioxide}{E}
\qa{Greenhouse gas from rice paddies, livestock, and landfills}{Methane}{A}
\qa{Natural warming by atmospheric gases}{Greenhouse effect}{E}
\qa{Graph of rising \ce{CO2} at Mauna Loa}{Keeling Curve}{D}
\qa{2015 pact to limit warming to well below \SI{2}{\celsius}}{Paris Agreement}{A}
\qa{UN scientific body assessing climate change}{IPCC}{A}
\qa{Treaty that phased out CFCs}{Montreal Protocol}{D}
\qa{Warming of the central--eastern equatorial Pacific}{El Ni\~no}{E}
\qa{El Ni\~no's effect on PH rainfall}{Below normal (drought)}{A}
\qa{La Ni\~na's effect on PH rainfall}{Above normal}{A}
\qa{ENSO stands for}{El Ni\~no--Southern Oscillation}{A}
\qa{Seawater becoming more acidic from absorbed \ce{CO2}}{Ocean acidification}{A}
\qa{Two causes of sea-level rise}{Thermal expansion + melting land ice}{D}
\qa{PH Climate Change Act}{RA 9729}{D}
\qa{PH Renewable Energy Act}{RA 9513}{D}
\qa{Wind farm in Ilocos Norte}{Bangui Wind Farm}{A}
\qa{Approximate global warming so far}{\verify{About \SI{1.1}{\celsius}}}{D}
\begin{confbox}
\confusion{Greenhouse effect}{ozone hole}{the greenhouse effect traps heat (\ce{CO2}); the ozone hole lets in more UV (CFCs). They are different problems.}
\confusion{El Ni\~no}{La Ni\~na}{in the PH, El Ni\~no means dry and La Ni\~na means wet.}
\confusion{Mitigation}{adaptation}{mitigation cuts the cause (emissions); adaptation copes with the effects (sea walls, drought-tolerant crops).}
\confusion{Global warming}{climate change}{warming is the temperature rise; climate change is the wider set of changes it drives.}
\end{confbox}

% ---------------------------------------------------------------
\subsection{Term 3: Plate Tectonics, Projectiles, Momentum, Electricity Supply}

\topic{Plate Tectonics: Drivers, Measurement, Mountain Building}
\kf
\begin{keyfacts}
\item Plate motion is measured with \textbf{GPS} (mm/yr precision), satellite methods, and seismic monitoring. Plates move a few cm per year, about as fast as fingernails grow.
\item Driving forces: mantle convection, \textbf{ridge push}, and \textbf{slab pull} (the strongest). The plates ride on the plastic asthenosphere.
\item Subduction of an oceanic plate makes a trench, a volcanic arc, and deep quakes along the Wadati--Benioff zone.
\item The \textbf{Himalayas} come from the India--Eurasia (continent--continent) collision. The \textbf{Andes} come from the Nazca Plate subducting under South America.
\item PH setting: the \textbf{Philippine Mobile Belt}, squeezed between the Philippine Sea Plate and the Eurasian (Sunda) Plate. Mt.\ Apo is the highest PH peak; the Sierra Madre is the longest PH range.
\item The Philippine Sea Plate moves at about \verify{\SIrange{7}{8}{cm} per year}. At \SI{8}{cm/yr} for 50 million years, that is about \SI{4000}{km}.
\end{keyfacts}
\cards
\qa{Satellite system that measures plate motion}{GPS}{E}
\qa{Strongest force driving plate motion}{Slab pull}{D}
\qa{Force from elevated mid-ocean ridges}{Ridge push}{A}
\qa{Heat-driven circulation in the mantle}{Mantle convection}{E}
\qa{Inclined zone of quakes along a subducting slab}{Wadati--Benioff zone}{D}
\qa{Plates whose collision built the Himalayas}{Indian and Eurasian}{A}
\qa{Plate subducting beneath South America}{Nazca Plate}{A}
\qa{Highest peak on Earth}{Mt.\ Everest}{E}
\qa{Highest mountain in the PH}{Mt.\ Apo}{E}
\qa{Longest mountain range in the PH}{Sierra Madre}{A}
\qa{Subducted plate melting forms}{A volcanic arc}{A}
\qa{Plate speeds are similar to}{Fingernail growth}{A}
\qa{Mobile zone the PH archipelago sits on}{Philippine Mobile Belt}{D}
\qa{\SI{8}{cm/yr} for 50 million years}{About \SI{4000}{km}}{D}
\begin{confbox}
\confusion{Ridge push}{slab pull}{ridge push shoves from the spreading ridge; slab pull drags from the sinking edge (stronger).}
\confusion{Andes}{Himalayas}{the Andes are ocean--continent (volcanic); the Himalayas are continent--continent (folded, few volcanoes).}
\end{confbox}

\topic{Projectile Motion}
\kf
\begin{keyfacts}
\item A projectile moves under gravity alone after launch, along a \textbf{parabolic} path.
\item With no air resistance, horizontal velocity is constant and vertical velocity changes by \SI{9.8}{\meter\per\second} every second.
\item At the peak, vertical velocity is 0 but the acceleration is still $g$ downward.
\item On level ground, maximum range is at \SI{45}{\degree}. Complementary angles (\SI{30}{\degree} and \SI{60}{\degree}) give the same range. Higher angles give a higher peak and longer flight.
\item A dropped ball and a ball thrown horizontally from the same height land at the same time.
\end{keyfacts}
\cards
\qa{Path of a projectile}{Parabola}{E}
\qa{Only force on an ideal projectile}{Gravity}{E}
\qa{Launch angle for maximum range}{\SI{45}{\degree}}{E}
\qa{Angle with the same range as \SI{30}{\degree}}{\SI{60}{\degree}}{A}
\qa{Horizontal velocity of an ideal projectile}{Constant}{A}
\qa{Vertical velocity at the peak}{Zero}{A}
\qa{Angle giving the greatest height}{\SI{90}{\degree}}{A}
\qa{Faster release at the same angle: range}{Increases}{E}
\qa{Acceleration at the peak}{$g$, downward}{D}
\qa{Dropped vs.\ horizontally thrown ball from the same height}{Land at the same time}{D}
\begin{confbox}
\confusion{Horizontal}{vertical motion}{horizontal is constant velocity; vertical is constant acceleration ($g$). They are independent.}
\confusion{Range}{maximum height}{\SI{45}{\degree} maximizes range; \SI{90}{\degree} maximizes height.}
\end{confbox}

\topic{Momentum and Collisions}
\kf
\begin{keyfacts}
\item Momentum $p = mv$, in \si{\kilogram\meter\per\second}. It is a vector. More momentum makes an object harder to stop.
\item \textbf{Impulse} $= F\Delta t = \Delta p$. Stretching the impact time lowers the force (seatbelts, airbags, crumple zones, helmets, bending your knees).
\item Total momentum before a collision equals total momentum after (closed system).
\item In an \textbf{elastic} collision, KE is conserved (bouncing). In an \textbf{inelastic} one, KE is lost. In a perfectly inelastic one, the objects stick. Momentum is conserved in all of them.
\end{keyfacts}
\cards
\qa{Formula for momentum}{$p = mv$}{E}
\qa{SI unit of momentum}{\si{\kilogram\meter\per\second}}{E}
\qa{Momentum of an object at rest}{Zero}{E}
\qa{\SI{50}{kg} runner at \SI{4}{\meter\per\second}}{\SI{200}{\kilogram\meter\per\second}}{A}
\qa{Force multiplied by time}{Impulse}{E}
\qa{Impulse equals the change in}{Momentum}{A}
\qa{Airbags reduce force by increasing}{Impact time}{A}
\qa{Car section designed to crush and absorb a crash}{Crumple zone}{E}
\qa{Loaded truck or bicycle at the same speed: harder to stop?}{Truck}{E}
\qa{Why bend your knees when landing}{Longer stopping time, less force}{A}
\qa{Collision in which KE is conserved}{Elastic}{E}
\qa{Objects stick together after colliding}{Perfectly inelastic}{A}
\qa{Is momentum conserved in inelastic collisions?}{Yes}{D}
\qa{Why a gun recoils}{Conservation of momentum}{A}
\qa{\SI{2}{kg} cart at \SI{3}{\meter\per\second} sticks to a resting \SI{1}{kg} cart: final speed}{\SI{2}{\meter\per\second}}{D}
\begin{confbox}
\confusion{Momentum}{inertia}{inertia depends on mass only; momentum needs motion ($mv$).}
\confusion{Elastic}{inelastic}{both conserve momentum; only elastic collisions also conserve kinetic energy.}
\confusion{Impulse}{force}{the same impulse spread over a longer time means a smaller force.}
\end{confbox}

\topic{Electricity: Generation, Transmission, Distribution}
\kf
\begin{keyfacts}
\item Generators turn mechanical energy into electrical energy by \textbf{electromagnetic induction} (Michael Faraday).
\item The chain: power plant $\rightarrow$ step-up transformer (high voltage means less loss) $\rightarrow$ transmission lines (NGCP) $\rightarrow$ substation step-down $\rightarrow$ distribution utility (Meralco or an electric cooperative) $\rightarrow$ meter $\rightarrow$ home.
\item Bills are in kilowatt-hours (kWh). The PH grid runs on AC at \SI{60}{Hz}.
\item Sources: renewable (geothermal, hydro, solar, wind, biomass) and non-renewable (coal, natural gas from Malampaya off Palawan, oil).
\end{keyfacts}
\cards
\qa{Principle behind generators}{Electromagnetic induction}{E}
\qa{Discovered electromagnetic induction}{Michael Faraday}{A}
\qa{Device that raises voltage for transmission}{Step-up transformer}{A}
\qa{Why power is transmitted at high voltage}{Less energy lost as heat}{D}
\qa{Operator of the PH transmission grid}{NGCP}{D}
\qa{Distribution utility of Metro Manila}{Meralco}{E}
\qa{Energy unit on electric bills}{Kilowatt-hour}{E}
\qa{\SI{1}{kW} appliance used for \SI{3}{h}}{\SI{3}{kWh}}{A}
\qa{Frequency of PH household AC}{\SI{60}{Hz}}{D}
\qa{Natural-gas field off Palawan}{Malampaya}{A}
\qa{Electricity from Earth's internal heat}{Geothermal}{E}
\begin{confbox}
\confusion{kW}{kWh}{kW is power (rate); kWh is energy (what you pay for).}
\confusion{Transmission}{distribution}{transmission is high-voltage, long-distance (NGCP); distribution is low-voltage, local (Meralco/co-ops).}
\confusion{AC}{DC}{AC reverses direction (wall outlets); DC flows one way (batteries).}
\end{confbox}

% ---------------------------------------------------------------
\subsection{Appendix: K-12 MELC-Only Grade 10 Topics}
{\footnotesize\textit{These topics are in the older Grade 10 MELCs but not in the MATATAG BOW. Study them only if your school or the contest still uses the K-12 MELCs.}\par}

\topic{Nervous System}
\kf
\begin{keyfacts}
\item A neuron has dendrites (receive), a cell body, and an axon (sends). The myelin sheath speeds the signal. Neurons meet at synapses, where neurotransmitters cross.
\item The CNS is the brain and spinal cord; the PNS is everything else. The autonomic system has sympathetic (``fight or flight'') and parasympathetic (``rest and digest'') branches.
\item The cerebrum handles thought and voluntary action, the cerebellum balance and coordination, and the medulla oblongata breathing and heartbeat.
\item Reflex arc: receptor $\rightarrow$ sensory neuron $\rightarrow$ interneuron $\rightarrow$ motor neuron $\rightarrow$ effector.
\end{keyfacts}
\cards
\qa{Branch of a neuron that receives signals}{Dendrite}{E}
\qa{Fatty layer that speeds nerve impulses}{Myelin sheath}{A}
\qa{Gap between two neurons}{Synapse}{A}
\qa{Brain part for balance}{Cerebellum}{E}
\qa{Brain part controlling heartbeat and breathing}{Medulla oblongata}{A}
\qa{``Fight or flight'' division}{Sympathetic}{D}

\topic{Endocrine and Reproductive Systems}
\kf
\begin{keyfacts}
\item The pituitary is the ``master gland.'' The thyroid makes thyroxine (metabolism; needs iodine). The adrenal glands make adrenaline.
\item Ovaries make estrogen and progesterone; testes make testosterone.
\item The menstrual cycle is about 28 days. FSH and LH come from the pituitary, and an LH surge triggers ovulation (around day 14). Fertilization usually happens in the fallopian tube.
\end{keyfacts}
\cards
\qa{Master gland}{Pituitary}{E}
\qa{Hormone that triggers ovulation}{LH}{D}
\qa{Usual site of fertilization}{Fallopian tube}{A}
\qa{Mineral needed to make thyroxine}{Iodine}{A}
\qa{Hormone that maintains the uterine lining}{Progesterone}{D}

\topic{Protein Synthesis (DNA \texorpdfstring{$\rightarrow$}{to} RNA \texorpdfstring{$\rightarrow$}{to} Protein)}
\kf
\begin{keyfacts}
\item The \textbf{central dogma} (Crick): DNA $\rightarrow$ RNA $\rightarrow$ protein.
\item \textbf{Transcription} copies DNA into mRNA in the nucleus. \textbf{Translation} builds the protein at the ribosome, with tRNA bringing amino acids.
\item A codon is 3 mRNA bases. The start codon is AUG (methionine). Stop codons are UAA, UAG, and UGA.
\end{keyfacts}
\cards
\qa{DNA $\rightarrow$ mRNA}{Transcription}{E}
\qa{mRNA $\rightarrow$ protein}{Translation}{E}
\qa{RNA that carries amino acids}{tRNA}{A}
\qa{Three-base unit on mRNA}{Codon}{A}
\qa{Start codon}{AUG}{D}

\topic{Gas Laws and Kinetic Molecular Theory}
\kf
\begin{keyfacts}
\item \textbf{Boyle}: $P \propto 1/V$ (constant $T$). \textbf{Charles}: $V \propto T$ (constant $P$, $T$ in K). \textbf{Gay-Lussac}: $P \propto T$ (constant $V$). \textbf{Avogadro}: $V \propto n$.
\item Ideal gas law: $PV = nRT$. At STP (\SI{273}{K}, \SI{1}{atm}), one mole of gas fills \SI{22.4}{L}. \SI{1}{atm} = \SI{760}{mmHg}.
\item Kinetic molecular theory: gas particles move randomly, have negligible volume and no attractions, and collide elastically. Average KE is proportional to the kelvin temperature.
\end{keyfacts}
\cards
\qa{Pressure up, volume down at constant $T$}{Boyle's law}{E}
\qa{Volume vs.\ temperature at constant $P$}{Charles's law}{E}
\qa{Pressure vs.\ temperature at constant $V$}{Gay-Lussac's law}{A}
\qa{Molar volume of a gas at STP}{\SI{22.4}{L}}{A}
\qa{\SI{1}{atm} in mmHg}{760}{A}
\qa{Gas at \SI{2}{L} and \SI{1}{atm} compressed to \SI{2}{atm} (constant $T$)}{\SI{1}{L}}{D}

\topic{Biomolecules}
\kf
\begin{keyfacts}
\item \textbf{Carbohydrates} (C, H, O) are the main energy source: glucose, sucrose, starch (plants), glycogen (animals), cellulose. Tests: Benedict's turns brick-red with reducing sugars; iodine turns blue-black with starch.
\item \textbf{Lipids} are fats, oils, phospholipids, and steroids, built from glycerol and fatty acids. They store long-term energy (grease-spot or Sudan III test).
\item \textbf{Proteins} are amino acids joined by peptide bonds; enzymes are proteins (Biuret test turns violet).
\item \textbf{Nucleic acids} are DNA and RNA, made of nucleotides.
\end{keyfacts}
\cards
\qa{Storage carbohydrate in animals}{Glycogen}{A}
\qa{Monomer of proteins}{Amino acid}{E}
\qa{Bond linking amino acids}{Peptide bond}{A}
\qa{Test for protein (turns violet)}{Biuret test}{D}
\qa{Benedict's test detects}{Reducing sugars}{D}

\topic{Electromagnetism: Motors and Generators}
\kf
\begin{keyfacts}
\item \textbf{\O{}rsted} (1820) showed that an electric current produces a magnetic field.
\item A motor turns electrical energy into mechanical energy; a generator does the reverse.
\item Fleming's left-hand rule is for motors; the right-hand rule is for generators.
\end{keyfacts}
\cards
\qa{Showed that current creates a magnetic field}{Hans Christian \O{}rsted}{A}
\qa{Device turning electrical into mechanical energy}{Electric motor}{E}
\qa{Fleming's rule for motors}{Left-hand rule}{D}
\begin{confbox}
\confusion{Motor}{generator}{a motor turns electricity into motion; a generator turns motion into electricity.}
\confusion{Boyle}{Charles}{Boyle's law is $P$--$V$ (inverse); Charles's law is $V$--$T$ (direct).}
\confusion{Transcription}{translation}{transcription writes DNA into RNA in the nucleus; translation reads RNA into protein at the ribosome.}
\end{confbox}

\end{gradebody}
\cardstats
